\PassOptionsToPackage{bookmarks=true,bookmarksopen=true,pdfborder={0 0 0},pdfhighlight={/N},linkbordercolor=[rgb]{0.5,0.5,0.5},implicit=false,hidelinks}{hyperref}
\usepackage[T1]{fontenc}   % so < > in \ttfamily render correctly, not as inverted marks
\usepackage{colortbl}      % \rowcolor in tabulars
\usepackage{cite}
\usepackage{amscd}
\usepackage{amsmath,amssymb,amsfonts}
\usepackage{cases}
\usepackage{commath}
\usepackage[hidelinks]{hyperref}
\usepackage{float}
\usepackage{forest}
\usepackage{hhline}
\usepackage{graphicx}

\usepackage{textcomp}
\usepackage{xcolor}
\usepackage{multirow}
\usepackage{tabularx,ragged2e}
\usepackage[skip=2pt]{caption}
\captionsetup[table]{skip=6pt}
\usepackage{subcaption}
\usepackage{multirow}
\usepackage{listings}
\usepackage{xcolor}
\usepackage{todonotes}
\usepackage{rotating}

\usepackage{booktabs} % For better table rules
\usepackage{threeparttable} % For table notes
\usepackage{multirow} % Required for merging rows
\usepackage{tabularx} % For adjustable-width columns
\usepackage{tabto}
\usepackage[most]{tcolorbox}
\usepackage{xcolor}
\usepackage{lipsum}
\usepackage{array}

% table spacing
\setlength{\tabcolsep}{8pt}  % column gap (default is 6pt)
\renewcommand{\arraystretch}{0.95}  % row height (default is 1.0)

\newcommand{\needsref}{\textsuperscript{*}}

\tcbset{
  user/.style={
    colback=blue!5!white,
    colframe=blue!75!black,
    fonttitle=\bfseries,
    title=User,
    sharp corners,
    boxrule=0.5pt,
    left=2mm, right=2mm, top=1mm, bottom=1mm,
  },
  agent/.style={
    colback=gray!10!white,
    colframe=black!60!white,
    fonttitle=\bfseries,
    title=Tool,
    sharp corners,
    boxrule=0.5pt,
    left=2mm, right=2mm, top=1mm, bottom=1mm,
  }
}

\newcommand\ipyt{\lstset{language=python}\lstinline}

\lstdefinestyle{matlab} {
        language=Matlab,
        backgroundcolor=\color[rgb]{0.95,0.95,0.92},
        commentstyle=\color[rgb]{0.13,0.55,0.13}\em,
        stringstyle=\color[rgb]{0.7,0,0},
        basicstyle=\ttfamily\scriptsize}

\lstset{style=matlab,
        captionpos=b,
        frame=single,
        keepspaces=true,
        breakautoindent=true, 
        breakindent=0pt, 
        breaklines=true,
        breakatwhitespace=true,
        numberbychapter=false
        }


\definecolor{backcolour}{rgb}{0.95,0.95,0.92}
\definecolor{codegreen}{rgb}{0,0.6,0}



% style for prompts in the papers sections
\lstdefinestyle{promptPaper}{
        language=Matlab,
        backgroundcolor=\color[rgb]{0.95,0.95,0.92},
        commentstyle=\color[rgb]{0.13,0.55,0.13}\em,
        stringstyle=\color[rgb]{0.7,0,0},
        basicstyle=\ttfamily\scriptsize
}

%json file style
\lstdefinelanguage{json}{
    basicstyle=\ttfamily\scriptsize,
    showstringspaces=false,
    breaklines=true,
    frame=single,
    backgroundcolor=\color[rgb]{0.95,0.95,0.92},
    string=[s]{"}{"},
    morecomment=[l]{:},
    morestring=[b]',
    literate=
     *{0}{{{\color{blue}0}}}{1}
      {1}{{{\color{blue}1}}}{1}
      {2}{{{\color{blue}2}}}{1}
      {3}{{{\color{blue}3}}}{1}
      {4}{{{\color{blue}4}}}{1}
      {5}{{{\color{blue}5}}}{1}
      {6}{{{\color{blue}6}}}{1}
      {7}{{{\color{blue}7}}}{1}
      {8}{{{\color{blue}8}}}{1}
      {9}{{{\color{blue}9}}}{1}
      {:}{{{\color{black}:}}}{1},
}

\def\arrow#1{\stackrel{#1}{\longrightarrow}}
\def\inp#1{^{\bullet}{#1}}
\def\out#1{{#1}^{\bullet}\!}
\def\exec#1{{[#1\rangle}}
\def\prioexec#1{{(#1\rangle}}
\def\free#1{\mathit{free} ( {#1} )}
\def\procs{\Pi}
\def\proc{\pi}
\def\subformula#1{\mathit{sub} (#1)}
\def\Pt{{\cal P}_t}
\def\fab{f_{a\leftrightarrow b}}
\def\avail{\mathit{avail}}
\def\F{\Diamond}
\def\G{\Box}


\def\cl{\mathit{cl}}
\def\arrow#1{\stackrel{#1}{\longrightarrow}}
%\def\best{\mathit{better}}
\def\true{\mathit{true}}
\def\bdd#1{\mathit{BDD}(#1)}
\def\btrue{\mathrm{BDD}(\top)}
\def\bfalse{\mathrm{BDD}(\bot)}
\def\max{\mathrm{BDD}(11\ldots11)}
\def\terms{\tt{terms}}
\def\counting{\mathit{count}}
\def\vars{\mathit{free}}
\def\proj{\mathit{hide}}
\def\ext{\mathit{ext}}
\def\domain{\mathit{domain}}
\def\extend{\mathit{extend}}
\def\lookup{\mathbf{lookup}}
\def\subf{\mathit{sf}}
\def\type{\mathit{type}}
\def\free{\mathit{free}}
\def\bddwedge{\mathbf{and}}
\def\BDDZ#1{\mathrm{BDD0}(#1)}
\def\BDDO#1{\mathrm{BDD1}(#1)}
\def\bddvee{\mathbf{or}}
\def\bddexists{\mathbf{exists}}
\def\bddset{\mathbf{build}}
\def\rename#1#2#3{\mathit{rename} (#1,\,#2,\,#3)}
\def\bddneg{\mathbf{not}}
\def\comp{\mathit{comp}}
\def\old{\mathsf{pre}}
\def\sem{{\cal S}}
\def\current{\mathsf{now}}
\def\nextold{\mathsf{nextpre}}
\def\false{\mathit{false}}
\def\pref#1#2{\mathit{pref}(#1,#2)}
\def\begini#1{\mathit{begin} ( #1 )}
\def\endi#1{\mathit{end} ( #1 )}
\def\maybe{\mathit{maybe}}
\def\pol{\mathit{pol}}
\def\xor{\oplus}
\def\tree#1{\mathit{tree}(#1)}
\def\asgn#1{\mathit{asgn}(#1)}
\def\var#1{\mathit{var}(#1)}
\def\lbl#1{\mathit{lbl}(#1)}
\def\since{\,{\cal S}\,}
\def\sincele#1{{\cal S}_{\leq #1}}
\def\esincele#1{{\cal Z}_{\leq #1}}
\def\sincel#1{{\cal S}_{< #1}}
\def\relle#1{{\cal R}_{\leq #1}}
\def\sinceg#1{{\cal S}_{> #1}}
\def\sincege#1{{\cal S}_{\ge #1}}
\def\relg#1{{\cal R}_{> #1}}
\def\add#1#2{{\mathit{add}(#1,#2)}}
\def\wsince{{\cal S}}
\def\empty{\mathit{empty}}
\def\released{{\cal P}}
\def\LTLsquare{\Box}
\def\BBox{\boxminus}
\def\LTLdiamond{\Diamond}
\def\LTLcircle{{\bigcirc}}
\def\GI{{\textfrak{I}}}
\def\SSI{{\textswab{I}}}
\newcommand\lola{\textsc{Lola}}
\newcommand\B{{\mathbb B}}
\newcommand\AP{\mathit{AP}}
\newcommand\buchi{B\"uchi}
\newcommand\cla{\mathsf{cone}}
\newcommand\sla{\mathsf{sla}}
\newcommand\suc{\mathsf{suc}}
\newcommand\clr{\mathsf{clr}}
\newcommand\cobuchi{Co-B\"uchi}
\newcommand\frag{\mathsf{frag}}
\newcommand\res{\mathsf{res}}
\newcommand\safe{\mathsf{safe}}
\newcommand\tempt{\mathsf{tempt}}
\newcommand\opt{\textrm{opt}}
\newcommand\best{\textrm{best}}
\newcommand\worst{\textrm{worst}}

%%%%%%%%%%%%%%%%
% Mathematics: %
%%%%%%%%%%%%%%%%

\newcommand{\boolty}{\mathbb{B}}

%%%%%%%%%%%%
% Symbols: %
%%%%%%%%%%%%

\newcommand\event[2]{{\em #1}(\text{#2})}
\newcommand\eventname[1]{{\em #1}}
\newcommand\args[1]{\text{#1}}
\newcommand{\symb}[1]{`$#1$'}
\newcommand{\seen}[1]{\mathit{}{seen}(#1)}
\newcommand{\alreadyseen}[1]{\mathit{}{\Theta}(#1)}

\newcommand\kwif{{\bf if\ }}
\newcommand\kwthen{{\bf then\ }}
\newcommand\kwelse{{\bf else\ }}

\newcommand\id[1]{\mathit{#1}}

\newcommand\DiamondMinus{\ensurestackMath{%
  \stackengine{.8pt}{\Diamond}{\scalebox{.65}[1.15]{$-$}}{2}{c}{F}{F}{L}}}

%%%%
\def\MnSymbolGlyphs#1{% IF ONE NEEDS TO LOCATE GLYPHS
  \usepackage{MnSymbol,fonttable}%
  \AtBeginDocument{\fonttable{MnSymbol#110}}%
}

\def\ImportFromMnSymbol#1{%
  \DeclareFontFamily{U} {MnSymbol#1}{}
  \DeclareFontShape{U}{MnSymbol#1}{m}{n}{
   <-6> MnSymbol#15
   <6-7> MnSymbol#16
   <7-8> MnSymbol#17
   <8-9> MnSymbol#18
   <9-10> MnSymbol#19
   <10-12> MnSymbol#110
   <12-> MnSymbol#112}{}
  \DeclareFontShape{U}{MnSymbol#1}{b}{n}{
   <-6> MnSymbol#1-Bold5
   <6-7> MnSymbol#1-Bold6
   <7-8> MnSymbol#1-Bold7
   <8-9> MnSymbol#1-Bold8
   <9-10> MnSymbol#1-Bold9
   <10-12> MnSymbol#1-Bold10
   <12-> MnSymbol#1-Bold12}{}
  \DeclareSymbolFont{MnSy#1} {U} {MnSymbol#1}{m}{n}
}
\newcommand\DeclareMnSymbol[4]{\DeclareMathSymbol{#1}{#2}{MnSy#3}{#4}}
% %%%%
\ImportFromMnSymbol{C}
\DeclareMnSymbol{\DDiamond}{\mathrel}{C}{120}

\renewcommand{\baselinestretch}{0.957}

% Klaus's macros

\usepackage{graphicx}
\usepackage{relsize}
\usepackage{stmaryrd}
\usepackage{algorithm}
\usepackage{algpseudocode}
\newcommand\matl{{\sf MaTL}}
\newcommand\para[1]{\vspace{2ex}\noindent{\bf #1.}}
\newcommand\subpara[1]{\underline{{\em #1.}}}
\newcommand{\Next}{\mathbin{\scalebox{0.8}{$\bigcirc$}}}
\newcommand{\prev}{\ominus}
\def\until{\;{\cal U}\;}
\def\wuntil{\;{\cal U_{\text{w}}}\;}
\def\wsince{\;{\cal S_{\text{w}}}\;}
\def\wnext{\,{\cal \Next_{\text{w}}}\,}
\def\wprev{\,{\cal \prev_{\text{w}}}\,}
\newcommand\ignore[1]{}
\renewcommand\true{\tt{tt}}
\renewcommand\false{\tt{ff}}
\renewcommand\iff{{\bf iff}}
\newcommand{\tabconst}{3.3cm}
\newcommand{\tabconste}{3.8cm}

\newcommand{\match}[3]{\mathsf{Match}\llbracket#1\rrbracket\left(#2\right)\left(#3\right)}
\newcommand{\pat}[3]{\mathsf{Pat}\llbracket#1\rrbracket\left(#2\right)\left(#3\right)}
\renewcommand{\eval}[2]{\mathsf{Eval}\llbracket#1\rrbracket\left(#2\right)}
\newcommand{\oper}[3]{\mathsf{Op}\llbracket#1\rrbracket\left(#2,\,#3\right)}
\newcommand{\rel}[3]{\mathsf{Rel}\llbracket#1\rrbracket\left(#2,\,#3\right)}
\newcommand{\reg}[2]{\mathsf{Reg}\llbracket#1\rrbracket\left(#2\right)}

\newcommand\nat{\mathbb{N}}
\newcommand\intg{\mathbb{Z}}
\newcommand\bool{\mathbb{B}}

\newcommand\fuzz{{\tt fuzz}}
\newcommand\sut{{\tt SUT}}

\definecolor{literalcolor}{rgb}{1,0,1}      % Pink for literals
\definecolor{operatorcolor}{rgb}{0,0.5,0.5} % Teal for operators (consistent with DejaVu)

\lstdefinelanguage{fuzzdsl}{
    keywords={
        rule, not, or, and, always, eventually, until, wuntil, next, wnext, sofar, once, since, wsince, prev, wprev, count, countpast, matches, any
    },
    literate=
        {|-}{{{\color{operatorcolor}$\vdash$}}}1
        {>}{{{\color{operatorcolor}$>$}}}1
        {<}{{{\color{operatorcolor}$<$}}}1
        {=}{{{\color{operatorcolor}$=$}}}2
        {>=}{{{\color{operatorcolor}$\ge$}}}1
        {<=}{{{\color{operatorcolor}$\le$}}}1
        {=>}{{{\color{red}$\Rightarrow$}}}1
        {[}{{{\color{red}$[$}}}1
        {]}{{{\color{red}$]$}}}1
        {<<}{{{\color{red}$<$}}}1
        {>>}{{{\color{red}$>$}}}1
        {+}{{{\color{operatorcolor}$+$}}}1
        {-}{{{\color{operatorcolor}$-$}}}1
        {*}{{{\color{operatorcolor}$*$}}}1
        {/}{{{\color{operatorcolor}$/$}}}1
        ,
    morecomment=[l]{\#},  
    backgroundcolor=\color{white},   
    commentstyle=\color{codegreen},
    keywordstyle=\color{blue},
    basicstyle=\ttfamily\scriptsize,
    breakatwhitespace=false,         
    breaklines=true,                 
    keepspaces=true,                      
    numbersep=5pt,                  
    showspaces=false,                
    showstringspaces=false,
    showtabs=false,                  
    tabsize=2,
    aboveskip=5pt,
    belowskip=2pt,
    frame=tblr,
    aboveskip=1em,    
    belowskip=1em,    
}

\newcommand{\fuzzdsl}{\lstset{language=fuzzdsl,backgroundcolor=\color{white}, frame=none}}
\newcommand{\ifuzz}{\fuzzdsl\lstinline}

\lstset{
  aboveskip=1em,    % space before each listing
  belowskip=1em,    % space after each listing
  basicstyle=\ttfamily\small,  % example style
  frame=single,     % optional: adds a frame around the listing
}

%\newcommand{\conjarrow}{\mathrel{\&\!\!\!\!\rightarrow}}

\lstset{
%  basicstyle=\ttfamily,
%  columns=fullflexible,
columns=fullflexible,keepspaces,
%    numberstyle=\tiny,
showstringspaces=false, % otherwise not possible to copy and paste
%        aboveskip=-40pt,
%        frame=leftline
% numbers=left,
  frame=single,
  language=Scala,
%  basicstyle=\small\ttfamily,
  basicstyle=\scriptsize\ttfamily,
  keywordstyle=\bfseries\color{red!60!blue},
  commentstyle=\ttfamily\color{tuggreen},
  stringstyle=\color{blue},
  backgroundcolor=\color{gray!10!white},
 morekeywords=[2]{1, 2, 3, 4, 5, 6, 7, 8, 9, 0},
  keywordstyle=[2]{\color{blue!20!pink}},
%  numberstyle=\tiny\color{gray},
%  emphlines={7},emphstyle=\color{tugred}
%  backgroundcolor=\color{yellow}
}

\lstdefinelanguage{tasknetdsl}{
    keywords={
        tasknet, 
        end, 
        timelines, 
        atomic,
        state,
        cumulative,
        bounds,
        rate,
        taskdef,
        pre,
        inv,
        maint,
        post,
        in,
        impacts,
        optional,
        task,
        durationrange,
        startrange,
        endrange,
        after,
        containedin,
        constraints,
        properties,
        prop,
        not,
        and,
        implies,
        or,
        always,
        eventually,
        until,
        sofar,
        once,
        since
    },
    literate=
        {|-}{{{\color{operatorcolor}$\vdash$}}}1
        {>}{{{\color{operatorcolor}$>$}}}1
        {<}{{{\color{operatorcolor}$<$}}}1
        {=}{{{\color{operatorcolor}$=$}}}2
        {>=}{{{\color{operatorcolor}$\ge$}}}1
        {<=}{{{\color{operatorcolor}$\le$}}}1
        {=>}{{{\color{red}$\Rightarrow$}}}1
        {[}{{{\color{red}$[$}}}1
        {]}{{{\color{red}$]$}}}1
        {<<}{{{\color{red}$<$}}}1
        {>>}{{{\color{red}$>$}}}1
        {+}{{{\color{operatorcolor}$+$}}}1
        {-}{{{\color{operatorcolor}$-$}}}1
        {*}{{{\color{operatorcolor}$*$}}}1
        {/}{{{\color{operatorcolor}$/$}}}1
        ,
    morecomment=[l]{\#},  
    backgroundcolor=\color{white},   
    commentstyle=\color{codegreen},
    keywordstyle=\color{blue},
    basicstyle=\ttfamily\scriptsize,
    breakatwhitespace=false,         
    breaklines=true,                 
    keepspaces=true,                      
    numbersep=5pt,                  
    showspaces=false,                
    showstringspaces=false,
    showtabs=false,                  
    tabsize=2,
    aboveskip=5pt,
    belowskip=2pt,
    frame=tblr,
    aboveskip=1em,    
    belowskip=1em,    
}

\newcommand{\tasknetdsl}{\lstset{language=tasknetdsl,backgroundcolor=\color{white}, frame=none}}
\newcommand{\itasknet}{\tasknetdsl\lstinline}

%%%%%%%%%%%%%%%%%%%%%%%%%%%%%%%%%%%%%%%%%%%%%%%%%%%%%%%%%%%%%%%%%%%%%%%%
% Promela (SPIN) listings language
%%%%%%%%%%%%%%%%%%%%%%%%%%%%%%%%%%%%%%%%%%%%%%%%%%%%%%%%%%%%%%%%%%%%%%%%

\lstdefinelanguage{promela}{
    morekeywords={
        active, assert, atomic, bit, bool, break, byte, chan, d_step,
        do, od, else, empty, enabled, eval, false, fi, full, goto, hidden,
        if, init, int, len, local, ltl, mtype, nempty, never, nfull, notrace,
        np_, od, of, pc_value, printf, priority, proctype, provided, run,
        short, show, skip, timeout, trace, true, typedef, unless, unsigned,
        xr, xs, inline, select, c_code, c_expr, c_decl, c_state, c_track
    },
    morekeywords=[2]{
        always, eventually, until, W, U, X
    },
    sensitive=true,
    morecomment=[l]{//},
    morecomment=[s]{/*}{*/},
    morestring=[b]",
    backgroundcolor=\color{gray!10!white},
    commentstyle=\color{codegreen}\itshape,
    keywordstyle=\bfseries\color{blue!70!black},
    keywordstyle=[2]{\bfseries\color{red!60!blue}},
    stringstyle=\color{brown},
    basicstyle=\ttfamily\scriptsize,
    columns=fullflexible,
    keepspaces=true,
    breaklines=true,
    breakatwhitespace=false,
    showspaces=false,
    showstringspaces=false,
    showtabs=false,
    tabsize=2,
    frame=single,
    aboveskip=0.6em,
    belowskip=0.6em,
    literate=
        {->}{{{\color{operatorcolor}$\rightarrow$}}}2
        {!}{{{\color{operatorcolor}!}}}1
        {?}{{{\color{operatorcolor}?}}}1
        {::}{{{\color{red}::}}}2
}

\newcommand{\promela}{\lstset{language=promela,backgroundcolor=\color{white},frame=none}}
\newcommand{\ipromela}{\promela\lstinline}

% Shell / terminal transcripts
\lstdefinestyle{shell}{
    basicstyle=\ttfamily\scriptsize,
    backgroundcolor=\color{black!5!white},
    frame=single,
    breaklines=true,
    columns=fullflexible,
    keepspaces=true,
    showstringspaces=false,
    commentstyle=\color{codegreen},
    aboveskip=0.6em,
    belowskip=0.6em,
}

% A procedurally generated starfield (no external image needed).
% Usage: \starfield  or  \starfield[0.6]  to scale.
\usepackage{tikz}
\newcommand{\starfield}[1][1.0]{%
\begin{tikzpicture}[scale=#1]
  \fill[black!90] (0,0) rectangle (4,2.5);
  \pgfmathsetseed{7}
  \foreach \i in {1,...,340}{%
    \pgfmathsetmacro{\sx}{rnd*4}%
    \pgfmathsetmacro{\sy}{rnd*2.5}%
    \pgfmathsetmacro{\sr}{0.004+rnd*0.016}%
    \pgfmathsetmacro{\so}{0.20+rnd*0.80}%
    \fill[white,opacity=\so] (\sx,\sy) circle (\sr);%
  }
  \foreach \i in {1,...,14}{%
    \pgfmathsetmacro{\bx}{rnd*4}%
    \pgfmathsetmacro{\by}{rnd*2.5}%
    \fill[white,opacity=0.15] (\bx,\by) circle (0.07);%
    \fill[white,opacity=0.35] (\bx,\by) circle (0.035);%
    \fill[white] (\bx,\by) circle (0.016);%
  }
  \draw[white!40,line width=0.3pt] (0,0) rectangle (4,2.5);
\end{tikzpicture}}

% Convenience macros for the SPIN talk
\newcommand{\spin}{{\sc Spin}}
\newcommand{\promelaname}{{\sc Promela}}
\newcommand{\tool}[1]{{\tt #1}}
\newcommand{\fprime}{F$^\prime$}